% 库拉托夫斯基闭包公理（Kuratowski closure axioms）（综述）
% license CCBYSA3
% type Wiki

本文根据 CC-BY-SA 协议转载翻译自维基百科\href{https://en.wikipedia.org/wiki/Kuratowski_closure_axioms}{相关文章}。

在拓扑学及其相关数学分支中，库拉托夫斯基闭包公理（Kuratowski closure axioms）是一组可用于在一个集合上定义拓扑结构的公理。它们与更常用的开集定义是等价的。该体系最早由卡齐米日·库拉托夫斯基（Kazimierz Kuratowski）形式化提出，\(^\text{[1]}\)后来又被数学家如瓦茨瓦夫·谢尔平斯基（Wacław Sierpiński）和安东尼奥·蒙泰罗（António Monteiro）等人进一步研究。\(^\text{[2]}\)

另一组类似的公理也可以仅通过其对偶概念——内点算子（interior operator）来定义拓扑结构。\(^\text{[3]}\)
\subsection{定义}
\subsubsection{库拉托夫斯基闭包算子及其弱化形式}
设 ( X ) 为任意集合，( \wp(X) ) 为其幂集。一个 **库拉托夫斯基闭包算子**（Kuratowski closure operator）是一个一元运算
[
\mathbf{c} : \wp(X) \to \wp(X)
]
满足以下性质：

* **[K1] 保持空集：**
  [
  \mathbf{c}(\varnothing) = \varnothing
  ]

* **[K2] 外延性（Extensive）：** 对所有 ( A \subseteq X )，
  [
  A \subseteq \mathbf{c}(A)
  ]

* **[K3] 幂等性（Idempotent）：** 对所有 ( A \subseteq X )，
  [
  \mathbf{c}(A) = \mathbf{c}(\mathbf{c}(A))
  ]

* **[K4] 对二元并集保持（或分配律）：** 对所有 ( A, B \subseteq X )，
  [
  \mathbf{c}(A \cup B) = \mathbf{c}(A) \cup \mathbf{c}(B)
  ]

由于闭包算子 (\mathbf{c}) 保持二元并集，上述条件的一个推论是：

* **[K4′] 单调性（Monotone）：**
  [
  A \subseteq B \Rightarrow \mathbf{c}(A) \subseteq \mathbf{c}(B)
  ]

事实上，如果我们将 [K4] 中的等号改为包含号，就得到一个更弱的公理 [K4″]（次可加性，subadditivity）：

* **[K4″] 次可加性：** 对所有 ( A, B \subseteq X )，
  [
  \mathbf{c}(A \cup B) \subseteq \mathbf{c}(A) \cup \mathbf{c}(B)
  ]

那么可以容易看出，公理 [K4′] 与 [K4″] 一起**等价于 [K4]**（见下文“证明2”的倒数第二段）。

库拉托夫斯基（1966）提出了一个**第五个（可选）公理**，要求单点集在闭包下保持不变：
对所有 ( x \in X )，
[
\mathbf{c}({x}) = {x}
]
他将满足全部五条公理的拓扑空间称为 **T₁ 空间（T1-spaces）**，以区别于仅满足前四条公理的一般空间。实际上，这些空间与通常意义下的拓扑 **T₁ 空间**完全对应（见下文）。[5]

如果省略条件 [K3]，则这些公理定义了一个 **Čech 闭包算子**。[6]
若改为省略 [K1]，则满足 [K2]、[K3] 与 [K4′] 的算子称为 **Moore 闭包算子**。[7]

因此，一个二元组
[
(X, \mathbf{c})
]
称为 **库拉托夫斯基空间（Kuratowski closure space）**、**Čech 闭包空间（Čech closure space）** 或 **Moore 闭包空间（Moore closure space）**，取决于闭包算子 (\mathbf{c}) 满足的公理体系。
